% -----------------------------------------------------------------------
\section{Introducción}
\label{sec:intro}

Los Problemas de Optimización Combinatoria (POC) aparecen en áreas tan
diversas como logística, manufactura, finanzas y planificación de recursos.
Su complejidad computacional, mayoritariamente \NP-dura, ha impulsado el
desarrollo de algoritmos exactos y, especialmente, de metaheurísticas capaces
de obtener soluciones de buena calidad en tiempos
prácticos~\cite{gendreau2010handbook,blum2003metaheuristics}.

La Búsqueda Tabú (TS)~\cite{glover1986future,glover1989tabu,glover1990tabu}
registra movimientos recientes en una lista tabú para evitar la re-visita
de soluciones y escapar de óptimos locales, y puede enriquecerse con
memoria de largo plazo, aspiración, reinicios y penalizaciones adaptativas.

La mayoría de los estudios compara TS contra otras familias
—algoritmos genéticos (GA)~\cite{holland1975adaptation}, recocido simulado
(SA)~\cite{kirkpatrick1983optimization} y optimización por colonias de
hormigas (ACO)~\cite{dorigo1996ant}— o propone variantes altamente
especializadas para un único
problema~\cite{cordeau2002guide,gendreau1994tabu,nowicki1996fast}.
Como advierte Hooker~\cite{hooker1995testing}, esa comparación competitiva
indica qué algoritmo gana, pero no qué componente explica el resultado.
Existe una brecha metodológica: \emph{repositorios públicos que ilustren
la progresión interna de TS} —de una configuración básica a una avanzada—
en múltiples problemas heterogéneos, con código reproducible y métricas
comparativas unificadas.

\subsection*{Contribuciones}

Las contribuciones del paper se resumen en:

\begin{enumerate}
  \item Implementación en Python de quince módulos de TS (5 problemas × 3
    configuraciones), coherentemente estructurados y documentados en español.
  \item Un orquestador de experimentos (\texttt{src/main.py}) que ejecuta los
    360 experimentos de forma automatizada, guarda resultados en JSON y
    genera tablas importables.
  \item Análisis cuantitativo del impacto de memoria de frecuencia, lista
    tabú dinámica y diversificación adaptativa en cinco POC clásicos.
  \item Un repositorio público que puede ser extendido a nuevos problemas o
    algoritmos sin modificar la infraestructura existente.
\end{enumerate}

La Sección~\ref{sec:marco} introduce los cinco problemas y los tres niveles
de TS.
La Sección~\ref{sec:metodologia} describe el diseño experimental y las
instancias de prueba.
La Sección~\ref{sec:resultados} presenta y analiza los resultados.
La Sección~\ref{sec:conclusiones} cierra con conclusiones y trabajo futuro.
